\section{Generalised Optics}
\label{sec:optics}

All optical organisation of the residual electromagnetic field is controlled by
the Angular Shortfall. Residual stress is the deflecting medium; no separate
optical axiom is required.

\subsection{Propagation zones}
\begin{center}
\begin{tabular}{lll}
\toprule
Zone & Overlap & Character \\
\midrule
Heart & 4-of-4 & Nearest vacuum; speed $\approx c$; minimal deflection \\
Skin & 3-of-4 & Intermediate residual; mild deflection \\
Membrane & 1--2 of 4 & Strongest residual; leakage, focusing, diffusion \\
Exterior & 0 & Ambient \\
\bottomrule
\end{tabular}
\end{center}

\subsection{Five optical modes}
\conceptual\ The residual field organises into five geometric regimes:

\begin{center}
\begin{tabular}{lll}
\toprule
Mode & Character & Geometric seat \\
\midrule
Caustic & Hyper-concentrated focus & Envelope of rays folded by shortfall gradients \\
Prismatic & Dispersion into components & Differential deflection at incomplete loci \\
Kinetic & Living intermediate residual & Ordinary residual motion; seat of $P/p$ \\
Scotopic & Dim, reduced intensity & Long membrane paths \\
Cryptic & Flat, outline-erasing diffusion & Solvent-boundary smoothing \\
\bottomrule
\end{tabular}
\end{center}

\subsection{Dual dissipation channels}
\begin{itemize}
\item \textbf{Caustic $=$ energetic}: local concentration, irreversible conversion
at singularity approach.
\item \textbf{Solvent boundary $=$ sympathetic}: resonant transfer across the
membrane between opposing residual domains.
\end{itemize}

\subsection{Refractive dipole}
Parallel to residual stress,
\[
N>1\quad(c_{\mathrm{med}}<c),\qquad n<1\quad(\text{phase }c_{\mathrm{med}}>c).
\]
Both are combinatorially forced once residual response is non-zero.
